\subsection{Problem 2: Dua Nakhoda, Satu Kapal}

\subsubsection{Identifikasi Masalah}
\begin{itemize}
\item \textbf{Kriteria 2 --- \textit{Leader election / mutual exclusion} dan \textit{effectively-once ingestion}:} tiga instance yang sama-sama melakukan \textit{polling} menggandakan permintaan ke upstream dan event di penyimpanan. Dibutuhkan tepat satu pemimpin aktif yang melakukan \textit{polling}, disertai penulisan idempoten agar pergantian pemimpin tidak menimbulkan duplikat maupun celah data.
\item \textbf{Kriteria 4 --- \textit{Unreliable clocks / clock skew} dan \textit{split-brain}:} jam fisik (\textit{wall-clock}) antar mesin bisa bergeser maju maupun mundur (+45 s / $-$30 s). Jika keabsahan \textit{lease} dihitung memakai jam lokal node Aggregator, dua instance bisa sama-sama merasa menjadi \textit{leader} (\textit{split-brain}) dan menulis bersamaan.
\item \textbf{Kriteria 5 --- \textit{Process pause / zombie leader} dan \textit{fencing token}:} \textit{leader} yang membeku (\textit{suspend} VM atau jeda GC panjang) melebihi batas waktu \textit{lease} tidak tahu bahwa kepemimpinannya sudah pindah ke instance lain. Ketika bangun kembali (\textit{zombie leader}), ia dapat menulis data basi ke penyimpanan. Penanganannya adalah \textit{fencing token} yang naik monoton dan diperiksa di sisi penyimpanan (Kleppmann, 2016. 2017, bab 8).
\item \textbf{Kriteria 6--8:} \textit{horizontal scaling} di bawah beban, \textit{load balancing} berbasis \textit{readiness}, dan \textit{rolling update} tanpa \textit{downtime}.
\end{itemize}

\subsubsection{Rancangan Solusi}
\begin{enumerate}
\item \textbf{Titik masuk tunggal dan \textit{graceful drain}.} Deployment \kode{aggregator} berjalan dengan 3 replika (\textit{soft podAntiAffinity} dan \kode{topologySpreadConstraints} per host) di belakang Service \kode{aggregator:8090}. Saat menerima \kode{SIGTERM}, \kode{GET /internal/ready} segera mengembalikan \kode{503 DRAINING}, worker dihentikan, \textit{lease} dilepas secara sukarela, dan pod menunggu \kode{preStop} / \kode{SHUTDOWN\_DRAIN\_S} 5 s agar \kode{kube-proxy} memutus trafik baru sebelum server HTTP ditutup. Strategi \textit{rolling update} memakai \kode{maxUnavailable: 0}, \kode{maxSurge: 1}.
\item \textbf{\textit{Lease} di CockroachDB (sumber kebenaran waktu = HLC \kode{now()}).} Tabel \kode{leader\_leases} diperbarui lewat \kode{AcquireOrRenewLease}. Seluruh perhitungan dan perbandingan kedaluwarsa (\kode{expires\_at > now()}) memakai \kode{now()} milik CockroachDB yang dijaga \textit{Hybrid Logical Clock} (HLC), tidak pernah memakai jam OS node Aggregator. Perpanjangan oleh pemegang sah mempertahankan \kode{fencing\_token}. Pengambilalihan setelah kedaluwarsa selalu menaikkan \kode{fencing\_token + 1}.
\item \textbf{\textit{Renew-before-poll}.} Sebelum mengirim permintaan ke BMKG/PVMBG pada setiap siklus, pemimpin memanggil \kode{elector.RenewNow()} secara sinkron ke DB. Jika \textit{zombie leader} baru bangun dari \textit{freeze}, panggilan ini langsung gagal dan instance mendemosi dirinya (\kode{leader\_lost}) tanpa menyentuh upstream.
\item \textbf{\textit{Fencing token} pada transaksi tulis (\kode{UpsertFencedBatch}).} Jika pemimpin membeku \textit{setelah} membaca upstream namun \textit{sebelum} menulis ke DB, \kode{UpsertFencedBatch} memverifikasi \kode{holder}, \kode{fencing\_token}, dan \kode{expires\_at > now()} di dalam transaksi \textit{serializable} yang sama dengan \textit{upsert} event, \textit{outbox}, dan pemajuan \textit{watermark}. Bila token sudah usang, transaksi di-\textit{rollback} (\kode{ErrFenced}, 0 baris tertulis) dan instance mendemosi dirinya.
\item \textbf{\textit{Watermark} durable dan idempotensi.} Tabel \kode{poller\_watermarks} menyimpan posisi terakhir secara atomik bersama batch tulis. Pemimpin baru memuat \textit{watermark} ini dengan mundur aman \kode{watermarkOverlap = 1 s} agar tidak ada event yang terlewat di batas detik, sementara kunci deterministik \kode{hazard\_id} (\kode{ON CONFLICT DO NOTHING/UPDATE}) mencegah duplikat. Mock mencatat permintaan per instance pada \kode{/admin/requests} untuk pembuktian.
\item \textbf{Uji beban (Fayadh).} k6 dijalankan dari luar klaster dengan skrip \kode{scripts/demo/m2/load/} (konfigurasi pada butir ``Di Laporan'' 4).
\end{enumerate}

\subsubsection{Alasan Keputusan}
\begin{itemize}
\item \textbf{\textit{Lease} di CockroachDB vs.\ Kubernetes \kode{Lease}:} \textit{fencing} hanya aman jika diperiksa tepat di tempat penulisan terjadi. Dengan menaruh \kode{leader\_leases} di CockroachDB yang sama, pengecekan \kode{fencing\_token} dan penulisan data berjalan atomik dalam satu transaksi \textit{serializable}, serta memakai satu sumber waktu (HLC CockroachDB). Kubernetes \kode{Lease} tidak bisa memagari transaksi SQL di basis data dan rentan terhadap pergeseran jam klien.
\item \textbf{Penetapan batas waktu dan \textit{trade-off}:} \kode{config.Load()} memvalidasi secara ketat bahwa \kode{LEASE\_TTL\_S > 2 * LEASE\_RENEW\_INTERVAL\_S}.
\begin{center}
\begin{tabularx}{\linewidth}{|p{4.3cm}|p{1.3cm}|X|}
\hline
\textbf{Parameter} & \textbf{Nilai} & \textbf{Fungsi dan \textit{trade-off}}\\ \hline
\kode{LEASE\_TTL\_S} & 15 s & Masa berlaku \textit{lease}. \textit{Worst-case failover} $\approx$ 15 + 5 = 20 s (jauh di bawah batas 60 s). Nilai terlalu kecil (mis.\ 5 s) memicu \textit{false failover} saat jeda GC/jaringan singkat. Nilai terlalu besar (mis.\ 45 s) mendesak \textit{failover} ke batas 60 s.\\ \hline
\kode{LEASE\_RENEW\_INTERVAL\_S} & 5 s & Jeda perpanjangan \textit{lease}. Memberi 3 kali kesempatan perpanjangan dalam 15 s sehingga gangguan satu permintaan tidak menjatuhkan \textit{leader}, sambil menjaga beban tulis sangat rendah (0,2 \textit{writes}/s).\\ \hline
\kode{safetyMargin} lokal & 2 s & Pemimpin lokal berhenti menganggap dirinya aktif pada 13 s tanpa perpanjangan sukses, sebagai rem konservatif sebelum \textit{lease} benar-benar habis di DB.\\ \hline
\kode{SHUTDOWN\_DRAIN\_S} & 5 s & Waktu tunggu setelah \kode{/internal/ready} mengembalikan 503 agar \textit{load balancer} menghentikan trafik baru sebelum proses keluar, tanpa melewati batas terminasi pod 30 s.\\ \hline
\kode{watermarkOverlap} & 1 s & Mundur 1 s saat \textit{polling} untuk mencegah celah di batas detik. Memicu pembacaan ulang beberapa event terbaru namun aman karena \textit{upsert} idempoten.\\ \hline
\end{tabularx}
\end{center}
\item \textbf{Uji beban dengan satu token per client:} 10.000 VU yang masing-masing meminta token akan menguji token-service, bukan jalur baca. Karena itu TTL token dinaikkan selama uji dan \textit{rate limit} per client dibuka (semua VU memakai satu client).
\end{itemize}

\subsubsection{Batasan yang Disadari}
\begin{itemize}
\item Keselamatan waktu bergantung pada batas \kode{--max-offset = 500 ms} CockroachDB. Jika jam OS pada sebuah node store melampaui 500 ms, node store tersebut mematikan dirinya sendiri untuk menjaga konsistensi HLC (mengurangi sementara kapasitas store).
\item Jika pemimpin membeku \textit{setelah} mengirim HTTP GET ke upstream tetapi \textit{sebelum} menulis ke DB, satu panggilan baca upstream tidak dapat dibatalkan. Namun \kode{UpsertFencedBatch} menjamin 0 baris basi ditulis ke basis data.
\item Uji beban 10.000 VU dari satu mesin dibatasi CPU, \textit{port}, dan jaringan mesin penguji. \isi{batasan yang teramati di klaster demo}
\end{itemize}

\subsubsection{Bukti Penyelesaian}

Skrip: \kode{scripts/demo/m2/p2/c1}--\kode{c5} (\kode{MODE=k8s}) dan \kode{scripts/demo/m2/load/}.

\begin{longtable}{|p{4.2cm}|p{10.7cm}|}
\hline
\textbf{Kriteria} & \textbf{Bukti}\\ \hline \endhead
1. 3 instance, satu titik masuk & \isi{\kode{c1-verification.txt}: instance yang di-\textit{drain} keluar dari \textit{endpoints}, 100\% HTTP 200}\\ \hline
2. Tepat sekali & \isi{\kode{c2}: \kode{duplicates: []} pada log mock, 0 duplikat \kode{(source, source\_ref\_id)}}\\ \hline
3. \textit{Failover} $\le$ 60 s & \isi{\kode{c3}: waktu pengambilalihan, \kode{fencing\_token} naik, 0 celah/duplikat (\kode{KILL\_MODE=node})}\\ \hline
4. \textit{Clock skew} +45/$-$30 s & \isi{\kode{c4-write-intervals.txt}: interval penulisan per \textit{epoch} tidak beririsan}\\ \hline
5. \textit{Freeze} & \isi{\kode{c5}: instance yang bangun tidak menulis (\kode{write rejected: stale fencing token})}\\ \hline
6. 10.000 VU, 3, 5, lalu 2 & \isi{\kode{evidence/load/<run>/verdict.txt}: total request, error tak terduga (\%), p95, error per fase. \kode{replicas.csv}}\\ \hline
7. \textit{Throughput} 3 vs 1 & \isi{\kode{throughput.sh}: served/s 1 replika, 3 replika, rasio}\\ \hline
8. \textit{Rolling update} & \isi{error fase \kode{rolling\_update} = 0}\\ \hline
\end{longtable}

\subsubsection{Jawaban atas Butir ``Di Laporan''}
\begin{enumerate}
\item \textbf{Persoalan:} kriteria 2 \textit{leader election / mutual exclusion} dan \textit{effectively-once ingestion}. Kriteria 4 \textit{unreliable clocks / clock skew} yang memicu \textit{split-brain}. Kriteria 5 \textit{process pause / zombie leader} yang dicegah dengan \textit{fencing token}.
\item \textbf{Mekanisme dan sumber waktu:} \textit{lease} di CockroachDB dengan \kode{now()} HLC sebagai satu-satunya otoritas waktu, dilengkapi \textit{renew-before-poll} dan pemeriksaan \kode{fencing\_token} di dalam transaksi \textit{serializable} penulisan. Pergeseran jam +45 s / $-$30 s pada node Aggregator (kriteria 4) tidak memengaruhi keputusan \textit{lease} karena jam lokal tidak pernah dibandingkan di DB. Instance yang bangun dari \textit{freeze} (kriteria 5) gagal pada \textit{renew-before-poll} atau ditolak oleh \kode{UpsertFencedBatch} (\kode{ErrFenced}) karena membawa \textit{token} lama, sehingga 0 baris basi tertulis.
\item \textbf{Batas waktu:} \textit{lease TTL} 15 s, perpanjangan 5 s, margin keselamatan lokal 2 s, \textit{shutdown drain} 5 s, dan \textit{watermark overlap} 1 s. \textit{worst-case failover} $\approx$ 20 s. Analisis \textit{trade-off} tiap nilai disajikan pada tabel bagian Alasan Keputusan di atas.
\item \textbf{Konfigurasi uji beban:}
\begin{center}
\begin{tabularx}{\textwidth}{|p{3.6cm}|X|}
\hline
Alat dan lokasi & k6 \isi{versi} di \isi{mesin di luar klaster}, menembak NodePort client-api (\kode{tests/load/k6\_m2.js})\\ \hline
Beban & 10.000 VU, \textit{ramp} 30 s, ditahan 150 s, jeda 0,5 s per iterasi, \kode{GET /v1/events?limit=20} dengan token Tim Lapangan\\ \hline
Jadwal & 0 s: 3 replika. 20 s: 5. 60 s: 2. 95 s: \textit{rolling update} (image baru atau revisi \textit{pod template})\\ \hline
Pengaturan selama uji & \textit{Rate limit} per client dibuka, TTL token 900 s (satu token per client dari \kode{setup()}), \textit{trace sampling} 1\%. Seluruhnya dikembalikan setelah uji\\ \hline
Klasifikasi & 200 dilayani. 429/503 penolakan terkontrol. Selain itu error. \textit{Threshold}: error $<$ 1\% dan 0 error pada fase \kode{rolling\_update}\\ \hline
\textit{Throughput} & 600 VU, 60 s per konfigurasi (1 lalu 3 replika), client-api 4 replika\\ \hline
\end{tabularx}
\end{center}
\end{enumerate}
